\section{Radial Collector Well}

% Describe source of problem
A radial collector well, also called a Ranney well, is a central caisson with several horizontal laterals that radiate outward near the base of an aquifer. The laterals give the well a large area of contact with the aquifer, so a radial collector well can yield much more water than a vertical well for the same drawdown. This example represents a radial collector well as a single multi-aquifer well in the \mf Multi-Aquifer Well (MAW) Package, with one vertical connection to the caisson and one horizontal connection to each cell a lateral passes through, using the NON\_VERTICAL\_WELLS option. A transient simulation compares the head in the well with the laterals represented as horizontal connections and as thin vertical screens.

\subsection{Example Description}

% spatial discretization and temporal discretization
The aquifer is 50 $m$ thick and is represented by one layer of 41 rows and 41 columns that are 25 $m$ wide. Horizontal and vertical hydraulic conductivity are 25 $m/d$ and specific storage is $1 \times 10^{-4}$ $1/m$. Head is 50 $m$ at the start of the simulation and is held at 50 $m$ in the cells around the edge of the domain. Model parameters are summarized in table~\ref{tab:ex-gwf-maw-collector-01}.

The caisson is in the center cell of the domain and four laterals extend 12 cells, 300 $m$, north, south, east, and west from it, so the well is connected to 49 cells: one vertical connection spanning the full thickness of the caisson cell and 48 lateral connections. The laterals are horizontal boreholes at mid-depth, 25 $m$ above the base of the aquifer. The well radius is 0.5 $m$ and a filter pack with a hydraulic conductivity of 25 $m/d$ extends to a radius of 1 $m$. The MEAN conductance equation is used for every connection, because the radial conductance equations assume that flow converges on the well in the horizontal plane, while flow to a horizontal lateral converges in a vertical plane.

With the NON\_VERTICAL\_WELLS option, each lateral connection is specified with a tilt angle of $90^{\circ}$ and a connection length of 25 $m$, the length of the lateral within the cell. The screen top and bottom are the top and bottom of the borehole, 1 $m$ apart, and determine the elevations over which the connection is saturated; the conductance of the connection depends on the 25 $m$ connection length. Without the option, the same connections are thin vertical screens 1 $m$ long, and the conductance of each lateral connection is calculated from that 1 $m$ screen length.

Two simulations are run. A steady-state simulation with the laterals represented as horizontal connections pumps 25,000 $m^3/d$ from the well. A transient simulation pumps 40,000 $m^3/d$ for 5 $d$ and then stops pumping for 5 $d$, each period divided into 25 time steps that increase by a factor of 1.2, and is run with the laterals represented as horizontal connections and as thin vertical screens. Storage in the well is not simulated in the transient simulation.

% add static parameter table(s)
\input{../tables/ex-gwf-maw-collector-01.tex}

% for examples without scenarios
\subsection{Example Results}

Simulated steady-state head is shown in figure~\ref{fig:ex-gwf-maw-collector-map}. The head in the well is 46.99 $m$, 3.01 $m$ below the starting head and within 0.004 $m$ of the head in the caisson cell. Drawdown extends along each lateral to its end, and the head contours are elongated along the laterals, because the well withdraws water from every cell a lateral passes through.

\begin{StandardFigure}{
    Simulated steady-state head for a radial collector well pumping 25,000 $m^3/d$, with the laterals represented as horizontal connections. Contours are labeled in meters. The open circle is the caisson, red cells are connected to the laterals, and gray cells have a constant head of 50 $m$.
    }{fig:ex-gwf-maw-collector-map}{../figures/ex-gwf-maw-collector-map.png}
\end{StandardFigure}

Simulated head in the well for the transient simulation is shown in figure~\ref{fig:ex-gwf-maw-collector-head}. With the laterals represented as horizontal connections, the head in the well at the end of pumping is 45.18 $m$, a drawdown of 4.82 $m$, and is 0.006 $m$ below the head in the caisson cell. With the laterals represented as thin vertical screens, the head in the well is 41.43 $m$, a drawdown of 8.57 $m$, and is 0.41 $m$ below the head in the caisson cell. The thin vertical screens have 1/25 of the conductance of the horizontal connections, so the laterals withdraw less water and more of the discharge enters the well through the caisson. At the end of pumping the head at the end of a lateral is 45.66 $m$ with horizontal connections and 46.34 $m$ with thin vertical screens, and the drawdown around the caisson is correspondingly larger with thin vertical screens. The head in the well recovers to within 0.1 $m$ of the starting head 0.93 $d$ after pumping stops with either representation of the laterals.

\begin{StandardFigure}{
    Simulated head in a radial collector well pumping 40,000 $m^3/d$ for 5 $d$ (shaded red) followed by 5 $d$ of recovery (shaded blue). The laterals are represented as horizontal connections, with a connection length equal to the length of the lateral in each cell, and as thin vertical screens one well diameter long.
    }{fig:ex-gwf-maw-collector-head}{../figures/ex-gwf-maw-collector-head.png}
\end{StandardFigure}
